% Appendix L — Reproducibility tiers
\section{Reproducibility}
\label{app:repro}

Each result row states which of five tiers a reader can reach with what we release
(\cref{tab:tiers}).

\begin{table}[h]
\caption{Reproduction tiers.}
\label{tab:tiers}
\centering\tablesf
\begin{tabularx}{\linewidth}{@{}llL@{}}
\toprule
Tier & Name & What a reader does \\
\midrule
R0 & read & checks receipts: hashes, counters, commands, exit codes and resource use \\
R1 & re-verify & re-runs the verifier on every committed artifact and recovers every utility \\
R2 & re-propose & runs the released checkpoint on the sealed states and recovers every committed artifact, byte for byte under the declared seed \\
R3 & retrain & trains a proposer from the released experience and recovers the primary results within their intervals \\
R4 & regenerate & regenerates experience with a teacher; protected for our teacher (\cref{app:disclosure}), open for any teacher the reader supplies \\
\bottomrule
\end{tabularx}
\end{table}

\paragraph{This draft} The draft itself is reproducible at R1 in the sense that matters for a
north-star document: every synthetic number and figure regenerates byte for byte from
\code{scripts/synthesize.sh} with a fixed seed, and \code{make check} refuses a draft in which any
synthetic value is unlabelled, any gap is unledgered or any figure plotting synthetic data is
untagged. \code{make submission-check} refuses this draft, by design.
